\cleardoublepage
\phantomsection
\begin{planbox}[แผนบริหารการสอนปฏิบัติการที่ 4\quad พลังงานกลและเครื่องปั่นเหวี่ยง Centrifuge (สัปดาห์ที่ 4)]
\noindent\textbf{1. เวลาเรียนปฏิบัติการ}\quad 3 ชั่วโมง (180 นาที)\quad \textbf{ผลลัพธ์การเรียนรู้ที่ประเมิน}\quad CLO1, CLO2, CLO3, CLO4

\noindent\textbf{2. วัตถุประสงค์เชิงพฤติกรรมประจำปฏิบัติการ}
\begin{enumerate}
    \item วัดรัศมีโรเตอร์ ($r_{\text{cm}}$) ของเครื่องปั่นเหวี่ยงและคำนวณแปลงค่าระหว่างความเร็วรอบต่อนาที ($\text{RPM}$) กับแรงเหวี่ยงสัมพัทธ์ ($\text{RCF}$ หรือ $\times g$) ได้ถูกต้อง \textbf{(CLO1, CLO2 เชื่อมโยง PLO2, PLO3)}
    \item ปฏิบัติการจัดสมดุลหลอดไมโครทิวบ์ในหัวปั่นโรเตอร์และปั่นแยกตะกอนเซลล์ยีสต์ออกจากส่วนใสได้อย่างปลอดภัย \textbf{(CLO2, CLO3, CLO4 เชื่อมโยง PLO1, PLO5, PLO6)}
\end{enumerate}

\noindent\textbf{3. แผนการจัดเวลาในคาบปฏิบัติการ 3 ชั่วโมง}
\begin{itemize}
    \item \textbf{ชั่วโมงที่ 1}\quad สาธิตโครงสร้างภายในของเครื่อง Centrifuge การวัดรัศมีโรเตอร์ และกฎเหล็กของการถ่วงสมดุลหลอดตรงข้ามกัน
    \item \textbf{ชั่วโมงที่ 2}\quad ปฏิบัติการวัดรัศมีโรเตอร์ คำนวณตารางแปลงค่า RPM $\leftrightarrow$ RCF และปั่นแยกตะกอนเซลล์ยีสต์ที่ความเร็วรอบต่างๆ
    \item \textbf{ชั่วโมงที่ 3}\quad เปรียบเทียบความแน่นของเม็ดตะกอนเซลล์ (Pellet) และความใสของส่วนใสด้านบน (Supernatant) พร้อมสรุปผล
\end{itemize}
\end{planbox}
\clearpage

\chapter{พลังงานกลและการคำนวณแรงเหวี่ยงสัมพัทธ์ของเครื่องปั่นเหวี่ยง}
\index{เครื่องปั่นเหวี่ยง Centrifuge}\index{แรงเหวี่ยงสัมพัทธ์ RCF}\index{เม็ดตะกอนเซลล์ Pellet}\index{การจัดสมดุลโรเตอร์}

\section{หลักการทางฟิสิกส์และทฤษฎีเบื้องหลังการทดลอง}

ในการเก็บเกี่ยวเซลล์แบคทีเรีย ยีสต์ หรือการตกตะกอนดีเอ็นเอในห้องปฏิบัติการจุลชีววิทยา เราใช้\textbf{เครื่องปั่นเหวี่ยง} (\textbf{Centrifuge}) เปลี่ยนพลังงานไฟฟ้าเป็นพลังงานจลน์ของการหมุนด้วยอัตราเร็วเชิงมุมสูง เพื่อให้เกิดความเร่งสู่ศูนย์กลางที่มีค่าสูงกว่าความเร่งโน้มถ่วงของโลก ($g$) หลายพันเท่า ค่าที่ระบุว่าแรงเหวี่ยงที่กระทำต่อเซลล์มีค่าเป็นกี่เท่าของแรงโน้มถ่วงโลก เรียกว่า \textbf{แรงเหวี่ยงสัมพัทธ์} (\textbf{Relative Centrifugal Force}, $\text{RCF}$ หรือ $\times g$)

\begin{conceptbox}[สมการแปลงความเร็วรอบ (RPM) เป็นแรงเหวี่ยงสัมพัทธ์ (RCF)]
\begin{equation}
    \text{RCF} = 1.118 \times 10^{-5} \times r_{\text{cm}} \times (\text{RPM})^2
\end{equation}
โดยที่ $r_{\text{cm}}$ คือรัศมีของโรเตอร์วัดจากจุดกึ่งกลางแกนหมุนไปถึงก้นหลอดตัวอย่างในหน่วยเซนติเมตร ($\text{cm}$) และ $\text{RPM}$ คือจำนวนรอบการหมุนต่อนาที
\end{conceptbox}

\begin{worked}[การคำนวณหาค่า RCF ของเครื่องปั่นเหวี่ยงไมโครทิวบ์]
\textbf{โจทย์ตัวอย่าง}\quad เครื่อง Microcentrifuge ในห้องแล็บจุลชีววิทยามีรัศมีโรเตอร์วัดถึงก้นหลอดได้ $r_{\text{cm}} = 8.5\text{ cm}$ หากตั้งความเร็วรอบที่ $6{,}000\text{ rpm}$ จะเกิดแรงเหวี่ยงสัมพัทธ์กี่ $\times g$

\textbf{วิธีคำนวณทีละขั้นตอน}
\[
    \text{RCF} = 1.118 \times 10^{-5} \times 8.5 \times (6{,}000)^2 = 1.118 \times 10^{-5} \times 8.5 \times (3.6 \times 10^7) = 3{,}421 \times g
\]
\end{worked}

\section{แผนภาพชุดอุปกรณ์การทดลองและข้อควรระวังความปลอดภัย}

\begin{figure}[H]
\centering
\resizebox{0.86\textwidth}{!}{%
\begin{tikzpicture}[font=\sffamily, >=Stealth]
    % Left: Measuring Rotor Radius r_cm
    \begin{scope}[shift={(0,0)}]
        \draw[fill=boxBlueBg, draw=physRoyalBlue, line width=1.2pt, rounded corners=4pt] (0,0) rectangle (6.8,4.2);
        \node[font=\bfseries\color{physRoyalBlue}] at (3.4,3.75) {(ก) การวัดรัศมีโรเตอร์ $r_{\text{cm}}$};
        \fill[physSlate!30, draw=charcoal, line width=1.3pt] (1.4,1.2) -- (5.4,1.2) -- (4.8,2.8) -- (2.0,2.8) -- cycle;
        \fill[charcoal] (3.25,0.8) rectangle (3.55,3.1);
        \draw[<->, crimsonAlert, line width=1.5pt] (3.4,1.5) -- node[above, font=\footnotesize\bfseries] {$r_{\text{cm}}$ ($\text{cm}$)} (5.1,1.5);
        \node[font=\scriptsize, align=center] at (3.4,0.4) {วัดจากกึ่งกลางแกนหมุนถึงก้นช่องใส่หลอด};
    \end{scope}

    % Right: Pellet vs Supernatant after Centrifugation
    \begin{scope}[shift={(7.6,0)}]
        \draw[fill=boxGreenBg, draw=bioEmerald, line width=1.2pt, rounded corners=4pt] (0,0) rectangle (6.8,4.2);
        \node[font=\bfseries\color{bioEmerald!90!black}] at (3.4,3.75) {(ข) หลอดไมโครทิวบ์ก่อนและหลังปั่นเหวี่ยง};
        % Before
        \fill[physAmber!45, draw=physRoyalBlue, line width=1.2pt] (1.3,1.2) -- (1.3,3.0) -- (2.3,3.0) -- (2.3,1.2) -- (1.8,0.6) -- cycle;
        \node[below, font=\scriptsize\bfseries] at (1.8,0.55) {ก่อนปั่น (ขุ่นทั่วหลอด)};
        \draw[->, physRoyalBlue, line width=1.5pt] (2.6,2.0) -- node[above, font=\scriptsize\bfseries] {$\text{RCF}$} (3.6,2.0);
        % After
        \fill[physCyan!15, draw=physRoyalBlue, line width=1.2pt] (4.1,1.2) -- (4.1,3.0) -- (5.1,3.0) -- (5.1,1.2) -- (4.6,0.6) -- cycle;
        \fill[physAmber!90!black] (4.35,0.9) -- (4.85,0.9) -- (4.6,0.6) -- cycle;
        \node[right, font=\scriptsize\bfseries, color=physRoyalBlue] at (5.15,2.3) {ส่วนใส (Supernatant)};
        \node[right, font=\scriptsize\bfseries, color=physAmber!90!black] at (4.95,0.8) {ตะกอนเซลล์ (Pellet)};
    \end{scope}
\end{tikzpicture}%
}
\caption{การวัดรัศมีโรเตอร์และการแยกชั้นระหว่างส่วนใสกับเม็ดตะกอนเซลล์หลังการปั่นเหวี่ยง}
\end{figure}

\begin{equipbox}[รายการอุปกรณ์สำหรับแต่ละกลุ่ม]
\begin{itemize}
    \item เครื่องปั่นเหวี่ยงหลอดไมโครทิวบ์ (Microcentrifuge), ไม้บรรทัดมิลลิเมตรสำหรับวัดรัศมีโรเตอร์ $r_{\text{cm}}$
    \item หลอดไมโครทิวบ์ขนาด $1.5\text{ mL}$ จำนวน 6 หลอด, ไมโครปิเปต P1000 และสารแขวนลอยเซลล์ยีสต์ \textit{Saccharomyces cerevisiae}
\end{itemize}
\end{equipbox}

\begin{warningbox}[กฎความปลอดภัยในการใช้เครื่อง Centrifuge]
ต้องวางหลอดที่มีปริมาตรเท่ากันเป๊ะในช่องตรงข้ามกัน ($180^\circ$) เสมอ ปิดฝาครอบในของโรเตอร์ให้แน่น และห้ามเปิดฝาเครื่องจนกว่าตัวเลขความเร็วรอบจะลดลงเป็น $0\text{ rpm}$
\end{warningbox}

\section{ขั้นตอนการปฏิบัติการทดลอง}
\begin{enumerate}
    \item ขณะเครื่องปั่นเหวี่ยงปิดสวิตช์อยู่ ให้ใช้ไม้บรรทัดวัดรัศมีโรเตอร์ ($r_{\text{cm}}$) จากกึ่งกลางแกนหมุนไปถึงขอบนอกของช่องใส่หลอด บันทึกค่าในหน่วยเซนติเมตร ($\text{cm}$)
    \item ใช้ไมโครปิเปต P1000 ดูดสารแขวนลอยยีสต์ที่เขย่าเข้ากันดีแล้ว ปริมาตร $1{,}000\text{ }\mu\text{L}$ ใส่ลงในหลอดไมโครทิวบ์ 3 หลอด และเตรียมหลอดน้ำกลั่นปริมาตร $1{,}000\text{ }\mu\text{L}$ อีก 3 หลอดสำหรับใช้เป็นหลอดถ่วงสมดุล (\textbf{Balance Tube})
    \item นำหลอดที่ 1 (พร้อมหลอดถ่วงสมดุลวางตรงข้ามกัน) ปั่นที่ความเร็วรอบ $1{,}500\text{ rpm}$ นาน $2\text{ นาที}$, หลอดที่ 2 ปั่นที่ $3{,}000\text{ rpm}$ นาน $2\text{ นาที}$ และหลอดที่ 3 ปั่นที่ $6{,}000\text{ rpm}$ นาน $2\text{ นาที}$ คำนวณค่า $\text{RCF}$ และสังเกตความใสของส่วนบนและขนาดของเม็ดตะกอนเซลล์ที่ก้นหลอด
\end{enumerate}

\section{ใบบันทึกผลการทดลองประจำกลุ่ม}

\begin{worksheetbox}[ใบบันทึกผลการทดลองที่ 4\quad การคำนวณค่า RCF และการปั่นเหวี่ยงแยกเซลล์จุลินทรีย์]
\noindent\textbf{กลุ่มปฏิบัติการที่} \dotfill \textbf{วันที่ทำการทดลอง} \dotfill \textbf{รัศมีโรเตอร์ที่วัดได้} $r_{\text{cm}} = $ \dotfill $\text{cm}$\\[4pt]
\noindent\textbf{รายชื่อสมาชิกกลุ่ม} 1. \dotfill 2. \dotfill 3. \dotfill

\vspace{6pt}
\noindent\textbf{ตารางบันทึกผลการคำนวณค่าแรงเหวี่ยงสัมพัทธ์ (RCF) และผลการตกตะกอนของเซลล์ยีสต์}
\begin{center}
\footnotesize\setlength{\tabcolsep}{3.5pt}
\begin{tabularx}{\linewidth}{c c c c >{\raggedright\arraybackslash}X}
\toprule
\textbf{หลอดที่} & \textbf{รอบ (RPM)} & \textbf{RCF คำนวณ ($\times g$)} & \textbf{RCF หน้าจอเครื่อง} & \textbf{ลักษณะส่วนใสและความแน่นของเม็ดตะกอน} \\
\midrule
หลอดที่ 1 & $1{,}500\text{ rpm}$ & ............ & ............ & \dotfill \\
หลอดที่ 2 & $3{,}000\text{ rpm}$ & ............ & ............ & \dotfill \\
หลอดที่ 3 & $6{,}000\text{ rpm}$ & ............ & ............ & \dotfill \\
\bottomrule
\end{tabularx}
\end{center}

\vspace{4pt}
\noindent\textbf{แสดงวิธีคำนวณค่า RCF ที่ความเร็วรอบ $3{,}000\text{ rpm}$ และ $6{,}000\text{ rpm}$}\\[4pt]
\dotfill\\[6pt]
\dotfill

\noindent\textbf{สรุปและอภิปรายผลการทดลอง}\\[4pt]
\dotfill\\[6pt]
\dotfill
\end{worksheetbox}

\section*{คำถามท้ายปฏิบัติการที่ 4}
\begin{enumerate}
    \item เมื่อเพิ่มความเร็วรอบ ($\text{RPM}$) จาก $3{,}000\text{ rpm}$ เป็น $6{,}000\text{ rpm}$ (เพิ่มขึ้น $2$ เท่า) ค่าแรงเหวี่ยงสัมพัทธ์ ($\text{RCF}$) เพิ่มขึ้นเป็นกี่เท่าของค่าเดิม เพราะเหตุใด
    \item เพราะเหตุใดหลังปั่นเหวี่ยงเสร็จแล้ว นักจุลชีววิทยาจึงต้องรีบดูดส่วนใส (Supernatant) ออกอย่างเบามือโดยไม่ให้ปลายทิปปิเปตไปชนก้นหลอดด้านที่หันออกนอกแกนหมุน
\end{enumerate}
